% chapters/notation.tex
% 符号表
% 源: 原书 PDF 第 587--588 页 (印刷页 583--584)

\chapter*{符号表}
\addcontentsline{toc}{chapter}{符号表}
\markboth{符号表}{符号表}

\section*{一般符号}

\begin{longtable}{@{}p{3.8cm}p{9.5cm}@{}}
$A^c$ & 集合 $A$ 的补集 \\
$E^\star$ & 对偶空间 \\
$\dual{\cdot}{\cdot}$ & 对偶 $E^\star,E$ 中的标量积 (对偶配对) \\
$[f=\alpha]=\{x;\ f(x)=\alpha\}$ & 泛函 $f$ 的等水平集 \\
$B(x_0,r)$ & 以 $x_0$ 为中心, $r$ 为半径的开球 \\
$B_E=\{x\in E;\ \norm{x}\le1\}$ & 空间 $E$ 中的闭单位球 \\
$\operatorname{epi}\varphi=\{[x,\lambda];\ \varphi(x)\le\lambda\}$ & 函数 $\varphi$ 的上图 \\
$\varphi^\star$ & 共轭函数 \\
$\LL(E,F)$ & 从 $E$ 到 $F$ 的有界线性算子空间 \\
$M^\perp$ & $M$ 的正交空间 \\
$D(A)$ & 算子 $A$ 的定义域 \\
$G(A)$ & 算子 $A$ 的图 \\
$N(A)$ & 算子 $A$ 的核 (零空间) \\
$R(A)$ & 算子 $A$ 的值域 \\
$\sigma(E,E^\star)$ & $E$ 上的弱拓扑 \\
$\sigma(E^\star,E)$ & $E^\star$ 上的弱$^\star$拓扑 \\
$\weakto$ & 弱收敛 \\
$J$ & 从 $E$ 到 $E^{\star\star}$ 的典范单射 \\
$p'$ & $p$ 的共轭指数, 即 $\dfrac{1}{p}+\dfrac{1}{p'}=1$ \\
几乎处处 (a.e.) & 在零测度集外 \\
$\abs{A}$ & 集合 $A$ 的测度 \\
$\operatorname{supp}f$ & 函数 $f$ 的支集 \\
$f\star g$ & $f$ 与 $g$ 的卷积 \\
$\rho_n$ & 磨光核序列 \\
$(\tau_hf)(x)=f(x+h)$ & 函数 $f$ 的平移 \\
$\omega\subset\subset\Omega$ & $\omega$ 紧含于 $\Omega$, 即 $\overline\omega$ 为紧集且 $\overline\omega\subset\Omega$ \\
$P_K$ & 到闭凸集 $K$ 上的投影算子 \\
$\abs{\cdot}$ & Hilbert 范数 \\
$\rho(T)$ & 算子 $T$ 的预解集 \\
$\sigma(T)$ & 算子 $T$ 的谱 \\
$\operatorname{EV}(T)$ & 算子 $T$ 的全体特征值的集合 \\
$J_\lambda=(I+\lambda A)^{-1}$ & 算子 $A$ 的预解式 \\
$A_\lambda=AJ_\lambda$ & 算子 $A$ 的 Yosida 逼近 \\
$\nabla u=\left(\dfrac{\partial u}{\partial x_1},\dfrac{\partial u}{\partial x_2},\dots,\dfrac{\partial u}{\partial x_N}\right)$ & 函数 $u$ 的梯度 \\
$D^\alpha u=\dfrac{\partial^{\abs{\alpha}}u}{\partial x_1^{\alpha_1}\partial x_2^{\alpha_2}\dots\partial x_N^{\alpha_N}}$ & 多重指标 $\alpha=(\alpha_1,\dots,\alpha_N)$, $\abs{\alpha}=\sum_{i=1}^N\alpha_i$ \\
$\Delta u=\sum_{i=1}^N\dfrac{\partial^2u}{\partial x_i^2}$ & $u$ 的 Laplace 算子 \\
$\R^N_+=\{x=(x',x_N)\in\R^{N-1}\times\R;\ x_N>0\}$ & 上半空间 \\
$Q=\{x=(x',x_N)\in\R^N\times\R;\ \abs{x'}<1\ \text{且}\ \abs{x_N}<1\}$ & 单位立方体 \\
$Q_+=Q\cap\R^N_+$ & 上半立方体 \\
$Q_0=\{x\in Q;\ x_N=0\}$ & 立方体的底面 \\
$(D_hu)(x)=\dfrac{1}{\abs{h}}(u(x+h)-u(x))$ & 差商算子 \\
$\dfrac{\partial u}{\partial n}$ & 外法向导数 \\
\end{longtable}

\section*{函数空间}

\begin{longtable}{@{}p{3.8cm}p{9.5cm}@{}}
$\Omega\subset\R^N$ & $\R^N$ 中的开集 \\
$\partial\Omega=\Gamma$ & $\Omega$ 的边界 \\
$L^p(\Omega)$ & $\{u:\Omega\to\R;\ u\ \text{可测且}\ \int_\Omega\abs{u}^p<\infty\}$, $1\le p<\infty$ \\
$L^\infty(\Omega)$ & $\{u:\Omega\to\R;\ u\ \text{可测且存在常数}\ C\ \text{使}\ \abs{u(x)}\le C\ \text{几乎处处}\}$ \\
$C_c(\Omega)$ & $\Omega$ 中具有紧支集的连续函数空间 \\
$C^k(\Omega)$ & $\Omega$ 上的 $k$ 次连续可微函数空间, $k\ge0$ \\
$C^\infty(\Omega)=\bigcap_{k\ge0}C^k(\Omega)$ & 光滑函数空间 \\
$C^k(\overline\Omega)$ & $C^k(\Omega)$ 中使得对满足 $\abs{\alpha}\le k$ 的每个多重指标 $\alpha$, 函数 $x\mapsto D^\alpha u(x)$ 可连续延拓到 $\overline\Omega$ 上的函数组成的集合 \\
$C^\infty(\overline\Omega)=\bigcap_{k\ge0}C^k(\overline\Omega)$ & 直到边界都光滑的函数空间 \\
$C^{0,\alpha}(\overline\Omega)$ & $\left\{u\in C(\overline\Omega);\ \sup_{\substack{x,y\in\Omega\\x\ne y}}\dfrac{\abs{u(x)-u(y)}}{\abs{x-y}^\alpha}<\infty\right\}$, $0<\alpha<1$ (Hölder 空间) \\
$C^{k,\alpha}(\overline\Omega)$ & $\{u\in C^k(\overline\Omega);\ D^ju\in C^{0,\alpha}(\overline\Omega)\ \forall j,\ \abs{j}\le k\}$ \\
$W^{1,p}(\Omega)$, $W_0^{1,p}(\Omega)$, $W^{m,p}(\Omega)$ & Sobolev 空间 \\
$H^1(\Omega)$, $H_0^1(\Omega)$, $H^m(\Omega)$ & 相应的 $L^2$ 型 Sobolev 空间 \\
\end{longtable}
